% !TeX root = ../main.tex

\chapter{Introduction}

Decentralized finance (DeFi) protocols encode financial logic directly in smart
contracts. A lending protocol determines whether a position is liquidatable from
collateral prices, an automated market maker updates liquidity and swap states
from pool values, and a vault accounts for deposits, debt, and rewards from
asset balances and reserve data. When these values are read from manipulable
spot prices, stale or misconfigured oracles, or insufficiently protected price
helper functions, attackers can temporarily distort the protocol's view of value
and induce incorrect economic behavior. As a result, price manipulation has
become an important and difficult vulnerability class in smart contract
security.

Flash-loan price manipulation is especially difficult to reason about because
the vulnerability is rarely a simple syntactic bug. The same price read may be
safe or unsafe depending on how it is computed, how fresh it is, what protection
surrounds it, and which financial decision later depends on it. For example,
using a current pool price directly in collateral valuation or liquidation logic
can be dangerous, while a properly configured time-weighted average price
(TWAP), independent oracle check, deviation bound, or slippage control may make
a similar-looking code pattern acceptable. Likewise, an accounting update,
token transfer, or reward calculation is only security-critical when the
manipulated value can influence an irreversible economic state transition. This
dependence on economic context makes price manipulation hard for detectors that
only match fixed code patterns.

Prior work has made substantial progress in detecting DeFi attacks and smart
contract vulnerabilities. Some systems analyze historical transactions to find
attacks that already occurred. Others inspect smart contract code or bytecode to
detect vulnerable patterns before exploitation. Recent approaches also combine
static analysis with large language models (LLMs) to reason about code context
and economic intent. These directions are useful, but their problem setting and
evaluation often remain centered on attack detection, individual contracts,
protocol-level vulnerability labels, or binary classification metrics.

Real audits have a different shape. Auditors typically inspect an entire
protocol repository rather than a single function or isolated transaction. They
must understand how contracts, helper libraries, price adapters, accounting
logic, and external integrations interact. The final audit report is also not a
binary label. A useful finding identifies the affected code location or snippet,
explains the root cause, describes the impact, recommends mitigation, and
provides enough evidence for the auditor and developer to verify the issue. A
tool that only says that a protocol is vulnerable may still leave most of the
audit work unresolved. A tool that reports many plausible but unsupported issues
may increase manual verification cost even if one of its reports is correct.

This thesis therefore studies flash-loan price manipulation detection as an
audit-assistance problem. The goal is not to replace human auditors, but to
help them quickly locate and understand suspicious price-dependent logic in a
complete repository. A useful assistant should gather concrete code evidence,
reduce redundant paths, explain why a particular price dependency is risky, and
produce findings whose structure resembles the information auditors already use:
code location, root cause, impact, missing protection, and supporting evidence.

To explore this goal, this thesis designs and implements a prototype workflow
for Solidity projects. The workflow uses Slither-based static analysis
\citep{slither} to collect candidate data-flow evidence related to price
dependencies and financial operations. It then groups redundant evidence into
representative entries and passes the reduced evidence to a staged LLM workflow.
The LLM stages are constrained to reason over the extracted evidence and to
emit structured findings rather than unconstrained prose. Each final finding is
designed to be traceable back to static-analysis artifacts, so that auditors can
inspect the evidence rather than trust a free-form model judgment.

This thesis also proposes an evaluation perspective aligned with audit
workflows. Protocol-level precision, recall, and F1-score are still useful as
baseline measurements, but they do not show whether a generated report points
to the right code, explains the right root cause, recommends the right
mitigation, or avoids excessive false positives. To evaluate audit usefulness,
this thesis uses public audit findings from sources such as Solodit, Code4rena,
and Sherlock as report-level comparison targets
\citep{solodit,code4rena,sherlock}. The evaluation considers localization,
semantic explanation similarity, impact and mitigation agreement, and
over-reporting penalties that approximate manual review burden.

In particular, this thesis is organized around three research questions. First,
can a static-analysis-plus-LLM workflow generate audit-style findings that help
auditors locate flash-loan price manipulation vulnerabilities in a repository?
Second, what failure modes are hidden when evaluation stops at binary
vulnerability classification? Third, can report-level evaluation based on audit
findings better reflect audit usefulness than protocol-level precision and
recall alone?

The contributions of this thesis are as follows:
\begin{enumerate}
    \item It frames flash-loan price manipulation detection as an
    audit-assistance problem, emphasizing repository-level localization,
    explanation quality, and false-positive burden rather than only vulnerable
    or benign classification.

    \item It presents a prototype workflow that combines static evidence
    extraction, path reduction, staged LLM reasoning, and structured
    audit-style finding generation for Solidity repositories.

    \item It proposes an audit-finding-based evaluation methodology that
    compares generated reports with public audit findings using localization,
    root-cause explanation, impact and mitigation alignment, semantic
    similarity, and over-reporting metrics.
\end{enumerate}

The remainder of this thesis is organized as follows. Chapter 2 reviews price
manipulation vulnerabilities, smart contract auditing practice, transaction- and
code-based detection approaches, LLM-assisted smart contract analysis, and
audit-report-based evaluation. Chapter 3 describes the design of the proposed
workflow and evaluation framework. Chapter 4 explains the implementation of the
prototype. Chapter 5 presents the planned evaluation methodology and
preliminary observations from existing artifacts. Chapter 6 concludes the
thesis and discusses future work.
